% =============================================================================
%  Lab 11 - Instruction Decoder II
% =============================================================================

\chapter{Lab 11 --- Instruction Decoder II}

\section{Objective}
% Two or three sentences, your own words.

\section{Specification --- the complete control table}

% THIS IS THE MOST IMPORTANT TABLE IN YOUR REPORT.
% Fill in every cell. Write "--" where CTRL_DEFAULT stands.
% PRINT THIS PAGE AND BRING IT TO LAB 12 - step three of the debugging method
% is "do the control signals match the table for this opcode?"

\begin{table}[htbp]
\centering
\caption{The complete control table. All eight instructions, all sixteen signals.
Differences from \texttt{CTRL\_DEFAULT} only.}
\label{tab:lab11-ctrl}
\scriptsize
\begin{tabular}{lllllllllll}
\toprule
instr & \texttt{ALU\_func} & \texttt{ALU\_Asel} & \texttt{ALU\_Ssel} & \texttt{IR\_Xsel} & \texttt{IR\_Xlsft} & \texttt{DM\_en} & \texttt{DM\_rw} & \texttt{DM\_strb} & \texttt{RF\_Dsel} & \texttt{RF\_strb} \\
\midrule
\texttt{ADD}  & & & & & & & & & & \\
\texttt{SLL}  & & & & & & & & & & \\
\texttt{ADDI} & & & & & & & & & & \\
\texttt{LUI}  & & & & & & & & & & \\
\texttt{LW}   & & & & & & & & & & \\
\texttt{SW}   & & & & & & & & & & \\
\texttt{BLEZ} & & & & & & & & & & \\
\texttt{J}    & & & & & & & & & & \\
\bottomrule
\end{tabular}

\vspace{6pt}

\begin{tabular}{lllll}
\toprule
instr & \texttt{RF\_Wsel} & \texttt{RF\_oe} & \texttt{NIAG\_sel} & \texttt{NIAG\_cond} \\
\midrule
\texttt{ADD}  & & & & \\
\texttt{SLL}  & & & & \\
\texttt{ADDI} & & & & \\
\texttt{LUI}  & & & & \\
\texttt{LW}   & & & & \\
\texttt{SW}   & & & & \\
\texttt{BLEZ} & & & & \\
\texttt{J}    & & & & \\
\bottomrule
\end{tabular}
\end{table}

\section{Design}

\subsection*{Three observations}
% REQUIRED, in your own words, one short paragraph each:
%
%   (a) ALU_Asel = IMMED for ADDI, LUI, LW and SW - four instructions with
%       nothing obvious in common. Why? What does that tell you about the cost
%       of a load compared with an add?
%
%   (b) DM_rw is determined by exactly ONE instruction out of eight, and most
%       of J's row is default. What does that say about the SHAPE of a control
%       table, and why does it make the default-first structure natural rather
%       than merely convenient?
%
%   (c) RF_strb is STROBE for exactly five instructions and HOLD for three.
%       State the rule that decides which.

\subsection*{Don't cares}
% REQUIRED. Example 2 gives two implementations of IR_Xsel that disagree on
% ADD, SLL and J. Explain:
%   - why both are correct
%   - which is cheaper
%   - and why a testbench that pinned IR_Xsel down for ADD would be testing an
%     accident of the implementation rather than a requirement.

\section{Implementation}

\vhdlfile{code/lab11/Decoder.vhd}{\texttt{Decoder} --- the complete instruction decoder}

\section{Verification}
% Your testbench is TABLE-DRIVEN: the expected control bundle for each
% instruction lives in an array, and one loop checks every field of every row.
% Say why that structure suits a table-shaped design, and report the count:
% "8 instructions, 14 fields each, N checks, all pass."
%
% Then explain the three latch tests and what class of bug they detect.

\begin{table}[htbp]
\centering
\caption{Latch tests --- the checks a one-instruction-at-a-time testbench
cannot make.}
\label{tab:lab11-latch}
\begin{tabular}{lll}
\toprule
sequence & signal checked & must be \\
\midrule
\texttt{LUI} after \texttt{ADDI} & \texttt{IR\_Xsel} & \\
\texttt{SW} after \texttt{BLEZ} & \texttt{IR\_Xlsft} & \\
\texttt{ADD} after \texttt{SW} & \texttt{RF\_oe} & \\
\bottomrule
\end{tabular}
\end{table}

\section{Results}
% Evidence: E1, E2, E4, E7.

\figplaceholder{the RTL Viewer for the complete \texttt{Decoder}}

% E1: 0 errors, and say EXPLICITLY that there is no inferred-latch warning.
% E4: register count MUST be zero.
% E7: quote the transcript line with the total check count.

\section{Discussion}
% REQUIRED. One bug, how you found it, what it cost.
%
% A good topic: did your expectation table and your decoder ever disagree?
% Which one turned out to be wrong - and what does that tell you about writing
% the test from the SPECIFICATION rather than from the implementation?
